\documentclass[final]{beamer}
\usepackage[size=custom,height=36in,width=48in,scale=1.0,orientation=landscape]{beamerposter}

% Color and font packages
\usepackage{graphicx}
\usepackage{booktabs}
\usepackage{multicol}
\usepackage{lmodern}
\usepackage{caption}
\usepackage{enumitem}
\usepackage[most]{tcolorbox}
\usepackage{amsmath,amssymb}

% Set colors
\definecolor{UUred}{RGB}{140,21,21}
\definecolor{UUgold}{RGB}{204,148,69}
\setbeamercolor{block title}{fg=white,bg=UUred}
\setbeamercolor{block body}{fg=black,bg=white}
\setbeamercolor{title in headline}{fg=UUred}

% Fonts
\setbeamerfont{block title}{size=\Large,series=\bfseries}
\setbeamerfont{block body}{size=\normalsize}
\setbeamerfont{title}{series=\bfseries,size=\LARGE}
\setbeamerfont{author}{size=\Large}
\renewcommand{\normalsize}{\fontsize{38}{48}\selectfont}
\renewcommand{\Large}{\fontsize{50}{52}\selectfont}
\renewcommand{\LARGE}{\fontsize{80}{52}\selectfont}
\setlist[itemize]{before=\fontsize{38}{48}\selectfont, label=\textbullet,
  topsep=2pt, itemsep=2pt, parsep=0pt, leftmargin=*}

% Tighter blocks
\setbeamertemplate{block begin}{%
  \begin{beamercolorbox}[ht=2.4em,colsep*=.5em]{block title}%
    \usebeamerfont*{block title}\insertblocktitle%
  \end{beamercolorbox}%
  \nointerlineskip%
  \begin{beamercolorbox}[colsep*=.75em,vmode]{block body}%
}
\setbeamertemplate{block end}{\end{beamercolorbox}\vskip0.5em}

% Poster info
\title{Addressing Competing Risks in High-Dimensional Survival Analysis: \\a Bayesian Spike-and-Slab LASSO Approach}
\author{Sophia Huebler, Xi Qiao, Boyi Guo}
\date{April 2026}
\institute{Population Health Sciences Department, University of Utah}

\begin{document}
\begin{frame}[t]

% ============================
% TITLE + LOGO HEADER
% ============================
\begin{columns}
  \begin{column}{0.10\textwidth}
    \includegraphics[height=8.5cm]{Designer.png}
  \end{column}
  \begin{column}{0.75\textwidth}
    \centering
    {\usebeamercolor[fg]{title in headline}
     \LARGE \textbf{\inserttitle}\par}
    \vspace{0.3em}
    \Large \insertauthor \\
    \Large \insertinstitute
  \end{column}
  \begin{column}{0.15\textwidth}
    \vspace{-2em}
    \includegraphics[width=0.95\linewidth]{UofU_Health_Logo.pdf}
  \end{column}
\end{columns}

\vspace{1cm}

% ============================
% MAIN BODY
% ============================
\begin{columns}[t, onlytextwidth]

% ============================
% LEFT COLUMN  (Context + SSL)
% ============================
\begin{column}{0.29\textwidth}

% --- Context ---
\begin{block}{Context}
\textbf{High-dimensional survival analysis can identify possible drug targets.}
\vspace{0.5cm}
\begin{columns}[c]
  \begin{column}{0.45\textwidth}
    \includegraphics[height=15cm]{Generic_Heatmap_Long.png}
  \end{column}
  \begin{column}{0.50\textwidth}
    \begin{itemize}
      \item $\#$ Samples $<\#$ Variables
      \item Small signals
      \item Meaningful correlation
      \item Time-to-event outcomes
      \item Censoring \& competing risks
    \end{itemize}
  \end{column}
\end{columns}

\vspace{0.6cm}
\textbf{Interpretability $>$ Predictive Accuracy}
\begin{itemize}
  \item Hypothesis generation: \emph{all} important variables matter
  \item Effect magnitude \& direction guide downstream prioritization
\end{itemize}

\vspace{0.6cm}
\textbf{Bayesian variable selection encodes sparsity through priors.}
\end{block}

% --- SSL Priors ---
\begin{block}{Spike-and-Slab LASSO Priors}
\textbf{Mixture of two double-exponential (DE) priors:} spike scale $s_0<$ slab scale $s_1$.
\begin{center}
  \includegraphics[height=12cm]{SSL_wide.png}
\end{center}
\begin{tcolorbox}[colback=red!5!white, colframe=red!75!black, title=Model Hierarchy]
\begin{align*}
\beta_j \mid \gamma_j, s_0, s_1 &\sim (1-\gamma_j)\,\text{DE}(0,s_0) + \gamma_j\,\text{DE}(0,s_1) \\
\gamma_j \mid \theta &\sim \text{Bernoulli}(\theta) \\
\theta &\sim \text{Beta}(1,1)
\end{align*}
\end{tcolorbox}
\end{block}

\end{column}

% ============================
% RIGHT COLUMN (Comp Risks + CIF/PSH/Wolber's)
% ============================
\begin{column}{0.70\textwidth}

% --- Competing Risks (top right) ---
\begin{block}{Competing Risks}
\begin{columns}[T]
  \begin{column}{0.40\textwidth}
    \textbf{A person who experiences a competing event cannot experience the event of interest.}
    \vspace{0.4cm}
    \begin{center}
      \includegraphics[height=14cm]{Comprisk.png}
    \end{center}
  \end{column}
  \begin{column}{0.58\textwidth}
    \begin{itemize}
      \item A variable can indirectly affect Event 1 by directly affecting Event 2
    \end{itemize}
    \vspace{0.3cm}
    \begin{center}
      \includegraphics[height=15cm]{placeholder.png}
    \end{center}
  \end{column}
\end{columns}
\end{block}

% ============================================================
% BOTTOM-RIGHT QUADRANT  --  the crux of the poster
% Linking model construction (PSH) and tuning (Wolber's C)
% through their shared target: the CIF
% ============================================================
\begin{block}{Linking Model and Tuning Through the Cumulative Incidence Function}

% --- CIF banner: the shared target ---
\begin{tcolorbox}[colback=UUgold!15!white, colframe=UUred!85!black,
                  boxrule=1.2mm, sharp corners,
                  left=4mm, right=4mm, top=2mm, bottom=2mm]
\begin{center}
\textbf{Shared target — the Cumulative Incidence Function:}\quad
$F_1(t \mid \mathbf{z}) = \Pr(T \le t,\; \epsilon = 1 \mid \mathbf{z})$
\end{center}
\vspace{-0.3em}
\begin{center}
\fontsize{34}{40}\selectfont
$\overbrace{\text{\textbf{PSH coefficients}}}^{\text{model structure}}
\;\Longrightarrow\;
\boxed{\;F_1(t\mid\mathbf{z})\;}
\;\Longleftarrow\;
\overbrace{\text{\textbf{Wolber's $C$-index}}}^{\text{model tuning}}$
\end{center}
\end{tcolorbox}

\vspace{0.4cm}

\begin{columns}[T]

% ----- Left half: PSH -----
\begin{column}{0.49\textwidth}
\begin{block}{Proportional Subdistribution Hazards}
\textbf{Fine \& Gray (JASA, 1999)}
\begin{tcolorbox}[colback=red!5!white, colframe=red!75!black,
                  title=Partial Pseudo Log-Likelihood]
With $N_i(t)=\mathbb{I}(T_i\le t,\epsilon_i=1)$ and at-risk $Y_i(t)=1-N_i(t-)$:
\vspace{-0.4em}
\begin{align*}
\ell_p(\boldsymbol\beta) = \sum_{i=1}^n \int_0^{\infty} \!\!&\!
  \Bigl[\mathbf{z}_i^\top\!\boldsymbol\beta -
   \log\!\bigl\{\!\textstyle\sum_q w_q(u)Y_q(u)e^{\mathbf{z}_q^\top\!\boldsymbol\beta}\bigr\}\Bigr] \\[-2pt]
  &\times w_i(u)\,dN_i(u)
\end{align*}
\vspace{-0.6em}

\noindent\textit{Time-dependent IPCW weights:}\;
$w_i(t)=\dfrac{\mathbb{I}(C_i \ge \min(T_i,t))\,\hat{G}(t)}{\hat{G}(\min(T_i,t))}$
\end{tcolorbox}
\begin{itemize}
  \item Subjects with the competing event \emph{stay} in the risk set
  \item Coefficients are \textbf{monotonically tied} to $F_1(t\mid\mathbf{z})$
\end{itemize}
\textbf{$\Rightarrow$ Model fits on the CIF.}
\end{block}
\end{column}

% ----- Right half: Wolber's C -----
\begin{column}{0.49\textwidth}
\begin{block}{Hyperparameter Tuning}
\textbf{Pre-validated prognostic scores} via cross-validation give an independent score for every subject. ``Risk'' of Event~1 is partially ordered:
\begin{itemize}
  \item $A < D < C$, \quad $B < C$, \quad $E < D < C$
\end{itemize}
\textbf{Wolber's $C$-index (Biostatistics, 2014)}
\begin{tcolorbox}[colback=red!5!white, colframe=red!75!black,
                  title=CIF-Aligned Concordance]
For prognostic model $M(t,\cdot)$ and pair $(i,j)$:
\vspace{-0.4em}
\begin{align*}
C_1(t) = \Pr\bigl(\,& M(t,X_i) > M(t,X_j) \,\bigm|\, \\
                   & \epsilon_i=1,\; T_i \le t,\; \\
                   & (T_j > T_i \text{ or } \epsilon_j = 2)\bigr)
\end{align*}
\end{tcolorbox}
\begin{itemize}
  \item Competing-event subjects are the \emph{lowest-risk} pairs
  \item Concordance is scored on the CIF's risk ordering
\end{itemize}
\textbf{$\Rightarrow$ Tuning evaluates on the CIF.}
\end{block}
\end{column}

\end{columns}

\vspace{0.3cm}

% --- Closing punchline banner ---
\begin{tcolorbox}[colback=UUred!90!black, colframe=UUred!90!black,
                  coltext=white, boxrule=0mm, sharp corners,
                  left=5mm, right=5mm, top=2mm, bottom=2mm]
\begin{center}
\textbf{Both model construction and tuning are anchored to $F_1(t\mid\mathbf{z})$ —
making PSH + Wolber's $C$ the aligned pair for high-dimensional competing-risks variable selection.}
\end{center}
\end{tcolorbox}

\end{block}

% --- Funding ---
\begin{block}{Funding and Acknowledgments}
Supported by NIH/NCI T32CA291646. The content is solely the responsibility of the authors and does not necessarily represent the official views of the funders.
\end{block}

\end{column} % right column
\end{columns}

\end{frame}
\end{document}
